\section{Zero mean dimension and regularity}
\label{sec:regularity-proof}

We now prove the characterization stated in \Cref{cor:regularity}.
The new equality supplies the implication from crossed-product
regularity back to zero mean dimension; the other directions use the
established results identified in the introduction.

\begin{proof}[Proof of \Cref{cor:regularity}]
Write $A=C(X)\rtimes_h\Z$ as in the corollary. This crossed product
is simple and nuclear
\cite[Section~3]{ElliottNiu2017}. It is also separable and unital, since
\(C(X)\) is separable and unital and the acting group is countable.
The canonical inclusion \(C(X)\subset A\) makes \(A\)
infinite-dimensional because \(X\) is infinite.

Lindenstrauss--Weiss proved that the small boundary property implies
zero mean dimension \cite[Theorem~5.4]{LindenstraussWeiss2000}.
Lindenstrauss proved the reverse implication for systems with an infinite
minimal factor \cite[Theorem~6.2]{Lindenstrauss1999}; here the identity
map supplies that factor. These directions are also recalled together
in \cite[Theorem~2.3]{ElliottNiu2017}. Elliott--Niu's Theorem~4.7 gives
Jiang--Su stability from zero mean dimension.

For the converse, suppose that \(A\) is Jiang--Su stable. Take
\(a,b\in M_\infty(A)_+\) with \(d_\tau(a)<d_\tau(b)\) for every
normalized trace \(\tau\) on \(A\), and place both elements in a common
finite corner \(M_q(A)\). This algebra is simple, exact, unital and
Jiang--Su stable. Every normalized trace on \(M_q(A)\) is
\(\tau\otimes\operatorname{tr}_q\) for a normalized trace \(\tau\) on
\(A\), where \(\operatorname{tr}_q=\operatorname{Tr}_q/q\). Thus, for
\(c\in M_q(A)_+\), its rank function for that normalized matrix trace is
\[
\lim_{m\to\infty}(\tau\otimes\operatorname{tr}_q)(c^{1/m})
   =\frac1q d_\tau(c),
\]
where the right side uses the unnormalized matrix extension fixed in
\Cref{eq:comparison-bound}. The strict rank inequalities therefore
hold for every normalized trace on \(M_q(A)\). R\o rdam's comparison
theorem \cite[Corollary~4.6]{Rordam2004} gives
\(a\precsim_{M_q(A)}b\), hence \(a\precsim b\) in the convention of this
paper. Thus \(0\) itself is a comparison bound, so \(\rc(A)=0\), and
\Cref{thm:main} gives
\(\mdim(X,h)=0\).

Finally, the crossed product satisfies all the hypotheses of
\cite[Theorems~A and~B]{CastillejosEvingtonTikuisisWhiteWinter2021}:
it is simple, separable, unital, nuclear and infinite-dimensional.
Finite nuclear dimension implies Jiang--Su stability by Winter's theorem
\cite{Winter2012}; the converse, with nuclear dimension at most one,
is supplied by Castillejos--Evington--Tikuisis--White--Winter.
Thus Jiang--Su stability, finite nuclear dimension and nuclear dimension
at most one are equivalent for $A$. This proves the five equivalences.
If $\mdim(X,h)>0$, including the value $\infty$, their contrapositives
give the final assertion.
\end{proof}
